\clearpage
\SourceLeaf{73}%
\pdfbookmark[1]{Göttingische gelehrte Anzeigen. 5. Dezember 1808.}{review1808}
\begin{center}\small Göttingische gelehrte Anzeigen. 5. Dezember 1808.\end{center}
\medskip
\enlargethispage{-9\baselineskip}
Das kürzlich erschienene Programm, in dem Professor \textsc{Gauß} seine Vorlesungen für das laufende Winterhalbjahr ankündigt, soll die Astronomen vor allem auf die Vorteile aufmerksam machen, die sich aus den heute zu Unrecht stark vernachlässigten \textit{Sternhöhen} für die Bestimmung von Zeit und Breite ergeben. Es lässt sich zwar nicht bestreiten, dass solche Beobachtungen besonders mit Reflexionsinstrumenten schwieriger sind und mehr Übung erfordern als Sonnenbeobachtungen. Ebenso gewiss ist aber, dass sich diese Schwierigkeiten überwinden lassen und sich die aufgewandte Mühe in sehr vielen Fällen vielfach auszahlt. Reisenden Beobachtern und Seeleuten ist die Verwendung von Sternhöhen besonders zu empfehlen. Sternhöhen können die oft vereitelten Sonnenhöhenmessungen vielfach ersetzen und bieten einen wichtigen Vorteil, der einem entgeht, wenn man sich auf Sonnenhöhen beschränkt. Aus einer einzelnen Sonnenhöhe kann man die Zeit nur bestimmen, wenn man die Breite bereits kennt. Umgekehrt muss die Zeit bekannt sein, wenn man die Breite aus einer Sonnenhöhe außerhalb des Meridians bestimmen will. Aus zwei Sonnenhöhen lassen sich Zeit und Breite nach bekannten Methoden durchaus gemeinsam bestimmen. Soll das Ergebnis aber einigermaßen genau sein, müssen diese Beobachtungen
\SourceLeaf{74}%
recht weit auseinanderliegen. Eine gute Uhr ist dabei unerlässlich, und auf See muss außerdem die Bewegung des Schiffes berücksichtigt werden. Verbindet man dagegen die beobachtete Höhe eines Sterns mit der Höhe eines zweiten Sterns, die man auf Wunsch unmittelbar danach messen kann, so lassen sich Zeit und Breite stets gemeinsam mit großer Genauigkeit bestimmen, sofern die Sterne selbst passend gewählt sind. Der Gang der Uhr muss dabei nur so zuverlässig bekannt sein, dass sich die Zeit zwischen den beiden Beobachtungen in Sternzeit umrechnen lässt. Die Aufgabe lässt sich ohne Weiteres auf die Lösung dreier sphärischer Dreiecke zurückführen. Dabei wird die numerische Rechnung allerdings etwas umfangreich. Professor \textsc{Gauß} hat jedoch auf rein analytischem Weg eine merklich kürzere Lösung entwickelt. Sie kann zugleich als Bestätigung dafür gelten, dass sich alles, was aus geometrischen Überlegungen folgt, mit analytischen Mitteln immer mindestens ebenso einfach und elegant finden lässt, wenn man sie \mbox{passend einsetzt.} Auch für den praktischen Gebrauch wird man die vorliegende Lösung weit bequemer finden als diejenige, die Herr \textsc{Kraft} im 13.~Band der \textit{Nova Acta Petropol.} für dasselbe Problem angegeben hat, nur mit der unwesentlichen Einschränkung, dass die beiden Höhenmessungen gleichzeitig erfolgen sollen. Der Verfasser verspricht, bei einer anderen Gelegenheit eine noch bequemere indirekte Lösung mitzuteilen. Die vorliegende Abhandlung erläutert außerdem, worauf man bei der Wahl der Sterne achten muss, damit die Ergebnisse möglichst genau ausfallen. Sie veranschaulicht das Verfahren durch seine tatsächliche Anwendung auf zwei an der hiesigen Sternwarte gemessene Sternhöhen. Die daraus abgeleitete Polhöhe stimmt bis auf $7''$ mit der von \textsc{Tob. Mayer} festgesetzten überein. Die daraus abgeleitete Zeit weicht nur um eine halbe Zeitsekunde von der Zeit ab, die am selben Tag auf andere Weise bestimmt wurde.
\par\bigskip\centerline{\rule{.24\textwidth}{.3pt}}
